\naslov{Kvazilinearne enačbe prvega reda v dveh spremenljivkah}

Obravnavamo enačbe oblike
\begin{equation}
  \label{eq:kle}
  a(x, y, u) u_x + b(x, y, u) u_y = c(x, y, u).
\end{equation}
Reševanja se bomo lotili z metodo karakteristik.
Enačbo prvo prepišemo v obliko
\[
  \sk{(a, b, c), (u_x, u_y, -1)} = 0.
\]
Ploskvi $z = u(x, y)$, ki reši kvazilinearno enačbo, pravimo \pojem{integralna
  ploskev} enačbe~\eqref{eq:kle}.
Če je $u: D \subseteq \R^2 \to \R$, lahko njen graf parametriziramo z
$\vec{r}(x, y) = (x, y, u(x, y))$.
Ploskev opišemo z normalo $\vec{r}_x \times \vec{r}_y = -(u_x, u_y, -1)$.
Če $u$ že imamo, tedaj vektor $(a, b, c)$ leži v tangentni ravnini na graf.

\begin{definicija}
  Vektor $(a, b, c)$ se imenuje \pojem{karakteristična smer} za PDE\@.
\end{definicija}

Pišimo $V = (a, b, c)$, čemur pravimo pripadajoče (karakteristično) vektorsko
polje.
Iz eksistenčnega izreka sledi, da za vsak $p_0 = (x_0, y_0, z_0)$ obstaja
natanko ena tokovnica polja $V$, ki poteka skozi točko $p_0$.
Če označimo $\varphi_t(p) = \gamma_p(t)$, je $\varphi$ tok, in velja
$\varphi_{t+s} = \varphi_t \circ \varphi_s$.

\begin{trditev}
  Naj bodo $a,b,c \in \zvodv{1}{\R^3}$, $S = \{z = u(x,y)\}$ poljubna integralna
  ploskev za~\eqref{eq:kle} in $p_0 \in S$.
  Tedaj vsaka karakteristična krivulja za~\eqref{eq:kle}, ki gre skozi točko
  $p_0$, cela leži v $S$.
\end{trditev}

\begin{proof}
  Naj bo $\gamma(t) = (x, y, z)$ karakteristična krivulja z $\gamma(0) = p_0$.
  Dokazati želimo, da velja $z = u(x, y)$, za kar je dovolj pokazati $\dot{z} =
  \partial_t u(x, y)$.
  Označimo $w(t) = z(t) - u(x(t), y(t))$ in računamo
  \begin{align*}
	\dot{w}
	&= \dot{z} - u_x \dot{x} - u_y \dot{y} \\
	&= c(\gamma(t)) - u_x(x, y) a(\gamma(t)) - u_y(x, y) b(\gamma(t)) \\
	&= c(x, y, w+u(x, y)) - u_x(x, y) a(x, y, w+u(x, y)) - u_y(x, y) b(x, y, w+u(x, y)).
  \end{align*}
  Velja $w(0) = 0$.
  Dobili smo Cauchyjevo nalogo oblike $\dot{w}(t) = f(t, w(t)), w(0) = 0$.
  Preverimo lahko, da za $f$ lokalno velja Lipschitzov pogoj, torej po
  eksistenčnem izreku obstaja enolična rešitev.
  Ker je $u$ po predpostavki rešitev~\eqref{eq:kle}, je $w = 0$ rešitev
  Cauchyjeve naloge, ki je enolična.
\end{proof}

\vprasanje{Pokaži, da tokovnice ne zapustijo integralne ploskve.}

\begin{posledica}
  Ob predpostavkah trditve je vsaka integralna ploskev unija karakterističnih krivulj.
\end{posledica}

\begin{posledica}
  Ob predpostavkah trditve:
  \begin{itemize}
  \item Če se dve integralski ploskvi $S_1, S_2$ sekata v točki $p \in S_1
	\presek S_2$, cela tokovnica skozi $p$ leži v $S_1 \presek S_2$.
  \item Če je $C = S_1 \presek S_2$ krivulja, ki je presek dveh integralskih
	ploskev $S_1, S_2$, ki se sekata ne-tangentno, je karakteristična krivulja.
  \end{itemize}
\end{posledica}

\begin{proof}
  Tokovnica ne zapusti ne $S_1$ ne $S_2$, torej mora biti vsebovana v preseku.
  Za drug del naj bo $p \in S_1 \presek S_2$.
  Tangentni ravnini $T_p S_1$ in $T_p S_2$ se sekata ne-tangentno, torej sta
  različni, in je njun presek enodimenzionalen.
  Ker sta $S_1$ in $S_2$ integralski ploskvi, njuni tangentni ravnini vsebujeta
  smer karakteristike in je zato
  \[
	T_p S_1 \presek T_p S_2 = \{\lambda(a(p), b(p), c(p)) \such \lambda \in \R\}
	= T_p(S_1 \presek S_2) = T_pC.
  \]
  Torej je $C$ v vsaki točki tangentna na smer karakteristike, torej je
  karakteristična krivulja.
\end{proof}

\vprasanje{Dokaži, da je presek integralskih ploskev karakteristična krivulja.}

Vsaka integralska ploskev $z = u(x, y)$ za~\eqref{eq:kle} je unija
karakterističnih krivulj polja $V = (a, b, c)$.
Naj bo $\Gamma$ gladka krivulja, parametrizirana z $\gamma(s) = (f(s), g(s),
h(s))$.
Skozi $\gamma(s)$ napeljemo karakteristično krivuljo, ki vsebuje točko
$\gamma(s)$.
Privzamemo lahko, da gre ta krivulja skozi točko $\gamma(s)$ ob času $t = 0$.
Iščemo vektorsko funkcijo $R(s, t) = (x(s, t), y(s, t), z(s, t))$, ki zadošča
\begin{itemize}
\item za vsak $s$ funkcija $R(s, \cdot)$ reši karakteristični sistem $\dot{x} =
  a(x, y, z), \dot{y} = b(x, y, z), \dot{z} = c(x, y, z)$,
\item začetni pogoji: $x(s, 0) = f(s), y(s, 0) = g(s), z(s, 0) = h(s)$.
\end{itemize}
Naj bo $J$ kompakten interval, $\gamma: J \to \R^3$, $s \in J$ in $\vec{F}_s =
(f(s), g(s), h(s))$.
Rešujemo sistem
\[
  \frac{d}{dt} \vec{x}_s = V(\vec{x}_s),\quad\vec{x}_s(0) = \vec{F}_s.
\]
Rešitev bo podana z $R(s, t) = \vec{x}_s(t)$.
Po eksistenčnem izreku za sisteme NDE obstaja $\varepsilon > 0$, da na intervalu
$(-\varepsilon, \varepsilon)$ obstaja natanko ena rešitev $\vec{x}_s:
(-\varepsilon, \varepsilon) \to \R^3$.
Želeli bi, da je $R$ parametrizacija integralske ploskve.
Problem je, če je $\Gamma$ karakteristična krivulja; tedaj je slika $R$
enodimenzionalna, torej potrebujemo dodatne pogoje.

Poskusimo sistem enačb $x = x(s, t), y = y(s, t)$ prevesti na sistem $s = s(x,
y), t = t(x, y)$ in vstaviti v tretjo komponento $R$.
To nam da $z = z(s(x, y), t(x, y)) = u(x, y)$.
Kdaj pa to smemo narediti?
Spomnimo se: če za neki $s_0 \in J$ velja
\[
  \frac{\partial (x, y)}{\partial (s, t)}(s_0, 0)
  =
  \begin{vmatrix}
	\frac{\partial x}{\partial s} & \frac{\partial x}{\partial t} \\
	\frac{\partial y}{\partial s} & \frac{\partial y}{\partial t}
  \end{vmatrix}
  \ne 0,
\]
inverz lahko poiščemo lokalno na neki okolici $(s_0, 0)$.
Po privzetkih na $\vec{F}_s$ je determinanta enaka
\[
  \begin{vmatrix}
	f'(s_0) & a(p_0) \\
	g'(s_0) & b(p_0)
  \end{vmatrix}
  =
  b(p_0) f'(s_0) - a(p_0) f'(s_0)
\]
za $p_0 = (f(s_0), g(s_0), h(s_0))$.
Geometrijsko želimo, da je projekcija na prvi dve komponenti tangente
$\dot{\gamma}(s_0)$ linearno neodvisna od projekcije vektorja $V(p_0)$ na prvi
dve komponenti.

\begin{trditev}
  Naj bo $\gamma = (f, g, h)$ pot v $\R^3$ razreda $\zvodvi{1}$, ki v dani točki
  $p_0 = \gamma(s_0)$ zadošča pogoju transverzalnosti
  \[
	\begin{vmatrix}
	  f'(s_0) & a(p_0) \\
	  g'(s_0) & b(p_0)
	\end{vmatrix}
	\ne 0.
  \]
  Tedaj v neki okolici točke $(x_0, y_0) = (f(s_0), g(s_0))$ obstaja natanko ena
  rešitev $u = u(x, y)$ kvazilinearne PDE~\eqref{eq:kle}, ki zadošča pogoju
  $h(s) = u(f(s), g(s))$ za vsak $s$ blizu $s_0$.
\end{trditev}

\begin{proof}
  Vse razen enoličnosti smo že pokazali.
  Denimo, da neka integralna ploskev vsebuje točko $p_0$.
  Tedaj po že dokazanem vsebuje vse karakteristične krivulje skozi točko $p_0$.
  Posledično (vsaj lokalno) celo ploskev parametriziramo z $(x(s, t), y(s, t),
  z(s, t))$.
  Če privzamemo, da v neki okolici $p_0$ obstaja še ena rešitev $w = w(x, y)$,
  za katero je $h = w(f, g)$ na okolici $s_0$, tedaj je $W = \{w(x, y) \such x,
  y\}$ unija tokovnic.
  Če gre neka tokovnica iz $W$ skozi točko $p_0$, je cela vsebovana v $U =
  \{u(x,y) \such x, y\}$.
  Torej $W \subseteq U$.
\end{proof}

\vprasanje{Kako poiščeš rešitev kvazilinearne PDE prvega reda dveh spremenljivk?
  Katere predpostavke potrebuješ? Pokaži, da je rešitev enolična.}

\podnaslov{Linearna PDE}

To je enačba oblike
\[
  a(x, y) u_x + b(x, y) u_y = c(x, y) u + d(x, y).
\]
V pripadajočem polju $V = (a, b, cu+d)$ funkcije $a, b, c, d$ niso odvisne od
$u$.
Enačbe karakteristik imajo obliko
\begin{align*}
  \dot{x} &= a(x, y), \\
  \dot{y} &= b(x, y), \\
  \dot{z} &= c(x, y) z + d(x, y).
\end{align*}

\vprasanje{Kakšne so enačbe karakteristik za linearno PDE?}

\podnaslov{Ovojnica družine ravnin}

Naj bo
\[
	\Pi_\lambda = \{(x, y, z) \in \R^3 \such z = \psi(x, y, \lambda) = a(\lambda)
	x + b(\lambda) y\}
\]
družina ravnin, kjer sta za interval $I \subseteq R$ funkciji $a, b \in
\zvodv{2}{I}$ taki, da je $\vec{n}(\lambda) = (a(\lambda), b(\lambda), -1)$
injektivna in
\[
	W(a', b') =
	\begin{vmatrix}
		a' & b' \\ a'' & b''
	\end{vmatrix}
	\ne 0
\]
za vsak $\lambda$.
\pojem{Ovojnica} družine $\Pi_\lambda$ je ploskev
\[
	\mathcal{O} = \{(x, y, z) \in \R^3 \such z = \varphi(x,y)\},
\]
za katero velja:
\begin{itemize}
\item $\mathcal{O} \subseteq \bigcup_\lambda \Pi_\lambda$,
\item če je $p \in \mathcal{O} \presek \Pi_\lambda$, potem je normala na
	$\mathcal{O}$ v točki $p$ vzporedna $\vec{n}(\lambda)$.
\end{itemize}
Trivialna izbira je $\mathcal{O} = \Pi_\lambda$ za nek $\lambda$, te pa v
nadaljnje ne bomo upoštevali.

\vprasanje{Definiraj ovojnico družine ravnin v $\R^3$.}

Vzemimo projekcijo $\pi : \R^3 \to \R^2$ na prvi dve komponenti.
Ker zahtevamo, da je $\vec{n}(\lambda)$ injektivna, za vsak par $(x, y) \ne
(0,0)$ obstaja natanko določena točka $p \in \mathcal{O}$, da je $(x, y) =
\pi(p)$.
To nam definira preslikavo $\Lambda : \pi(\mathcal{O}) \to \R$, ki slika par
$(x, y)$ v pripadajoči $\lambda$.

Naš cilj je poiskati funkcijo $\varphi$, ki definira $\mathcal{O}$.
Ker je $\mathcal{O} \subseteq \bigcup_\lambda \Pi_\lambda$, velja
\[
	\varphi(x, y) = a(\Lambda(x, y)) x + b(\Lambda(x, y)) y.
\]
Zahtevali smo, da je normala vzporedna z $\vec{n}(\lambda)$, torej lahko
parametriziramo $\mathcal{O}$ z
\[
	u(x, y) = (x, y, \varphi(x, y)).
\]
Upoštevaje definicijo $\Pi_\lambda$ lahko zapišemo normalo kot
\[
	u_x \times u_y
	= (1, 0, \varphi_x) \times (0, 1, \varphi_y)
	= (- \psi_x - \psi_\lambda \Lambda_x, -\phi_y - \psi_\lambda \Lambda_y, 1)
\]
oziroma nasprotno vrednost kot
\[
	-u_x \times u_y = (\psi_x, \psi_y, -1) + \psi_\lambda (\Lambda_x, \Lambda_y,
	0).
\]
Velja $\psi_x = a$ in $\psi_y = b$, torej $-u_x \times u_y = \vec{n}(\Lambda(x,
y)) + \psi_\lambda(\Lambda_x, \Lambda_y, 0)$.
Zahtevamo, da je to vzporedno z $\vec{n}(\Lambda(x, y))$, torej mora biti drugi
člen ničeln.
Če je $\Lambda_x = \Lambda_y = 0$, je $\Lambda$ konstantna preslikava in je
$\mathcal{O} = \Pi_\Lambda$; to je izključen trivialni primer.
Velja torej $\psi_\lambda = 0$, kar nam da dodatno enačbo za ovojnico.

\vprasanje{Kako poiščeš ovojnico družine ploskev v $\R^3$? Izpelji potrebni
	pogoj.}

Ta enačba je ekvivalentna pogoju
\[
	D(\lambda) := \frac{a'(\lambda)}{b'(\lambda)} = \frac{-y}{x}.
\]
Po predpostavki o determinanti Wronskega je odvod $D$ neničeln, torej jo lahko
na množici $\{x \ne 0\}$ obrnemo in dobimo $\Lambda_1(x, y) =
D^{-1}(\nicefrac{-y}{x})$.
Podobno lahko naredimo za $y \ne 0$, iz česar izpeljemo $\Lambda_2(x, y) =
E^{-1}(\nicefrac{-x}{y})$ za $E(\lambda) = \nicefrac{b'(\lambda)}{a'(\lambda)}$.
Če pokažemo, da je $\Lambda_1 = \Lambda_2$, kjer sta obe definirani, bomo lahko
zapisali
\[
	\Lambda(x, y) =
	\begin{cases}
		\Lambda_1(x, y), & x \ne 0 \\
		\Lambda_2(x, y). & y \ne 0
	\end{cases}
\]
Vemo, da za $(x, y) \in \pi(\mathcal{O})$ obstaja natanko določen $\lambda$, pri
katerem je $(x, y)$ slika neke točke na $\mathcal{O}$ s $\pi$.
Ker je tretja koordinata $z = a(\Lambda_1) x + b(\Lambda_1) y = a(\Lambda_2) x +
b(\Lambda_2) y$ pri tem natančno določena, iz pogoja injektivnosti $\vec{n}$
lahko izpeljemo $\Lambda_1 = \Lambda_2$.

Iz zgornje izpeljave hitro vidimo, da velja $\Lambda(\sigma x, \sigma y) =
\Lambda(x, y)$ za $\sigma \ne 0$, iz česar sledi naslednja trditev.

\begin{trditev}
	Če je $p \in \mathcal{O}$, potem je $\sigma p \in \mathcal{O}$ za poljuben
	$\sigma \ne 0$.
\end{trditev}

% LocalWords:  privzetkih transverzalnosti Wronskega
